\section*{第 \ref{ch:b1:optical-mineralogy} 章——光学矿物学}

\begin{solution}{exo:b1:optical-mineralogy:1}
石英：$\qty{30000}{nm}\times0.009 = \qty{270}{nm}$，一级白色。
方解石：$\qty{30000}{nm}\times0.172 = \qty{5160}{nm}$，远超色谱表范围：
淡的高级白色，略带粉红或绿色。
\end{solution}

\begin{solution}{exo:b1:optical-mineralogy:2}
$I/I_{\max} = \sin^2 2\theta$：在 \ang{0} 处为 $0$；$\sin^2\ang{45} = 0.5$ 对应
\ang{22.5}；在 \ang{45} 处为 $1$；在 \ang{67.5} 处为 $0.5$；在 \ang{90} 处为 $0$。
\end{solution}

\begin{solution}{exo:b1:optical-mineralogy:3}
\omterm{def:b1:rock-forming-minerals:groups}{石榴子石}属立方晶系，因而是\omterm{def:b1:optical-mineralogy:anisotropic}{均质矿物}：光穿过时保持原有的振动方向，
即起偏镜方向，被检偏镜挡住；在每个位置均有 $\Delta = 0$。
垂直于\omterm{def:b1:optical-mineralogy:optic-axis}{光轴}切割的石英颗粒，沿该轴传播的光同样满足
$n_1 = n_2$，故 $\Delta = 0$，同样保持暗场。
在\omterm{def:b1:field-lab-methods-1:microscope}{单偏光}下，\omterm{def:b1:rock-forming-minerals:groups}{石榴子石}的\omterm{def:b1:optical-mineralogy:relief}{突起}很高（$1.83$，
封片剂为 $1.54$），石英的\omterm{def:b1:optical-mineralogy:relief}{突起}几乎为零。
\end{solution}

\begin{solution}{exo:b1:optical-mineralogy:4}
$|n - 1.54|$：铁铝榴石 $0.29$，镁\omterm{def:b1:rock-forming-minerals:groups}{橄榄石} $0.11$，正\omterm{def:b1:rock-forming-minerals:groups}{长石} $0.02$，石英
$0.01$；\omterm{def:b1:optical-mineralogy:relief}{突起}按此顺序由高到低。正\omterm{def:b1:rock-forming-minerals:groups}{长石}折射率较低，
为负\omterm{def:b1:optical-mineralogy:relief}{突起}，升高焦点时其\omterm{def:b1:optical-mineralogy:relief}{贝克线}向颗粒外移动，进入
封片剂一侧。
\end{solution}

\begin{solution}{exo:b1:optical-mineralogy:5}
$\delta = 650/30000 = 0.022$。一种绿色的\omterm{def:b1:optical-mineralogy:pleochroism}{多色性}\omterm{def:b1:rock-forming-minerals:mineral}{矿物}，其解理夹角为
\ang{56} 和 \ang{124}，\omterm{def:b1:optical-mineralogy:birefringence}{双折射率}约为 $0.020$：\omterm{def:b1:rock-forming-minerals:groups}{角闪石}。
\end{solution}

\begin{solution}{exo:b1:optical-mineralogy:6}
$\tan\ang{30} = 0.5774$。石英：$(0.013/1.54)\times0.5774 = \qty{0.00487}{rad}
= \ang{0.28}$。镁\omterm{def:b1:rock-forming-minerals:groups}{橄榄石}：$(0.111/1.54)\times0.5774 = \qty{0.0416}{rad} =
\ang{2.38}$。比值 $0.111/0.013 = 8.5$。精确计算：$\sin r = 1.54\sin\ang{30}/1.651
= 0.4664$，$r = \ang{27.80}$，偏折角为 \ang{2.20}；折射率差如此之大，
一级近似公式高估了 \qty{8}{\percent}。
\end{solution}

\begin{solution}{exo:b1:optical-mineralogy:7}
$270 + 540 = \qty{810}{nm}$，为二级绿色，此时颗粒与试板的慢\omterm{def:b1:optical-mineralogy:optic-axis}{光轴}方向平行，
即延迟相加（延迟相加）；$|270 - 540| =
\qty{270}{nm}$，重新回到一级白色，此时颗粒的慢\omterm{def:b1:optical-mineralogy:optic-axis}{光轴}
与试板的快\omterm{def:b1:optical-mineralogy:optic-axis}{光轴}方向平行。
\end{solution}

\begin{solution}{exo:b1:optical-mineralogy:8}
全厚度处 $\Delta = \qty{30000}{nm}\times0.052 = \qty{1560}{nm}$。从
边缘（$\Delta = 0$）向内，\omterm{def:b1:optical-mineralogy:retardation}{光程差}依次经过 542 和
\qty{1074}{nm}，但未经过 \qty{1605}{nm}：两条红紫色条纹。
\end{solution}

\begin{solution}{exo:b1:optical-mineralogy:9}
在颗粒上测得的\omterm{def:b1:optical-mineralogy:extinction}{消光角}取决于颗粒的切割方向：
只有特定方向的切面才显示最大\omterm{def:b1:optical-mineralogy:extinction}{消光角}，其他切面显示的角度较小，
直至为零。应比较的是最大值，约为 \ang{37}，
它与透\omterm{def:b1:rock-forming-minerals:groups}{辉石}（约 \ang{38}）相符，而非\omterm{def:b1:rock-forming-minerals:groups}{角闪石}（至多
\ang{19}）。
\end{solution}

\begin{solution}{exo:b1:optical-mineralogy:10}
检偏镜沿 $x$ 方向时，两束波的 $x$ 分量为
$E_0\cos\theta\cdot\cos\theta$ 与 $-E_0\sin\theta\,e^{-i\varphi}\cdot(-\sin\theta)$，
故 $A = E_0(\cos^2\theta + \sin^2\theta\,e^{-i\varphi})$，且
\[
  |A|^2 = E_0^2\bigl(\cos^4\theta + \sin^4\theta + 2\sin^2\theta\cos^2\theta\cos\varphi\bigr)
  = E_0^2\bigl(1 - 2\sin^2\theta\cos^2\theta\,(1 - \cos\varphi)\bigr),
\]
利用 $\cos^4\theta + \sin^4\theta = 1 - 2\sin^2\theta\cos^2\theta$。由于
$2\sin^2\theta\cos^2\theta = \frac12\sin^2 2\theta$ 与 $1 - \cos\varphi =
2\sin^2(\varphi/2)$, $I_\parallel = I_0(1 - \sin^2 2\theta\sin^2(\pi\Delta/\lambda))$.
对每一波长，$I_\parallel + I_\perp = I_0$：平行偏光下的颜色
与\omterm{def:b1:field-lab-methods-1:microscope}{正交偏光}下的颜色互为补色。一颗\omterm{def:b1:optical-mineralogy:retardation}{光程差}为
\qty{550}{nm} 的颗粒，在\omterm{def:b1:field-lab-methods-1:microscope}{正交偏光}下呈品红色，
在平行偏光下则为淡黄绿色。
\end{solution}

\begin{solution}{exo:b1:optical-mineralogy:11}
(a) $\Delta/\lambda$ 从 $5000/700 = 7.14$ 变化到 $5000/400 = 12.5$；它
经过整数 8 至 12：共五个零点。(b) 零点位于
$\lambda_k = \Delta/k$；相邻两个零点之间的距离为 $\Delta/k - \Delta/(k + 1) =
\Delta/(k(k + 1)) \approx \lambda^2/\Delta$，在 \qty{550}{nm} 附近即 $550^2/5000 =
\qty{60.5}{nm}$。(c) 光谱中分布的五条窄暗带会吸收每种颜色的一小部分，
其余光以平均二分之一的权重通过，
人眼看到的是淡雅、近乎中性的色调。
\end{solution}

\begin{solution}{exo:b1:optical-mineralogy:12}
$\sin^2 2\theta$ 在一周内的平均值为 $\frac1{2\pi}\int_0^{2\pi}\sin^2 2\theta\,
\mathrm d\theta = \frac12$，因此平均强度为
$\frac12 I_0\sin^2(\pi\Delta/\lambda)$。她计算时取 $\delta = \Delta/30000$
而非 $\Delta/25000$（单位为 nm）：她的所有数值均为真实值的 $25/30 = 0.83$，
偏低 \qty{17}{\percent}。
\end{solution}

\begin{solution}{pb:b1:optical-mineralogy:1}
\textbf{1.} B：暗边表示\omterm{def:b1:optical-mineralogy:relief}{突起}高；A 的边界几乎看不见，说明
\omterm{def:b1:optical-mineralogy:relief}{突起}低。

\textbf{2.} A：$(1.544 + 1.553)/2 = 1.5485$，$|n - n_m| = 0.0085$。B：
$(1.664 + 1.672 + 1.694)/3 = 1.6767$，$|n - n_m| = 0.1367$。

\textbf{3.} 偏折量与 $|n - n_m|$ 成正比：$0.1367/0.0085 =
16.1$。

\textbf{4.} B 的折射率高于封片剂：正\omterm{def:b1:optical-mineralogy:relief}{突起}。

\textbf{5.} \omterm{def:b1:optical-mineralogy:pleochroism}{多色性}；黑\omterm{def:b1:rock-forming-minerals:groups}{云母}。

\textbf{6.} 两组解理位于 \ang{56} 与 \ang{124}：为\omterm{def:b1:rock-forming-minerals:groups}{角闪石}。

\textbf{7.} $\qty{30000}{nm}\times0.009 = \qty{270}{nm}$：一级白。

\textbf{8.} 各颗粒的切面方向不同，所显示的
$|n_1 - n_2|$ 介于 0 与 $\delta$ 之间；只有平行于\omterm{def:b1:optical-mineralogy:optic-axis}{光轴}切割的
颗粒才显示 $\delta$ 本身。在众多颗粒中，
应取\omterm{def:b1:optical-mineralogy:interference-colour}{干涉色}最高者。

\textbf{9.} $d = 342/0.009 = \qty{38000}{nm} = \qty{38}{\micro m}$.

\textbf{10.} $(38 - 30)/30 = \qty{26.7}{\percent}$.

\textbf{11.} $\sin^2(\pi\times342/550) = 0.861$；
当厚度为 \qty{30}{\micro m} 时，$\sin^2(\pi\times270/550) = 0.999$。

\textbf{12.} $\qty{38000}{nm}\times0.008 = \qty{304}{nm}$ 至
$\qty{38000}{nm}\times0.013 = \qty{494}{nm}$：由白色至橙红色，均低于
一级的上限 \qty{542}{nm}。

\textbf{13.} 介于 1074 与 \qty{1605}{nm} 之间：三级。

\textbf{14.} $\delta = 1140/38000 = 0.030$.

\textbf{15.} $1140/30000 = 0.038$，这提示可能是\omterm{def:b1:rock-forming-minerals:groups}{橄榄石}。A
镁\omterm{def:b1:rock-forming-minerals:groups}{橄榄石}与石英反应生成\omterm{def:b1:rock-forming-minerals:groups}{辉石}，
\ce{Mg2SiO4 + SiO2 -> 2 MgSiO3}；一种与石英同时结晶的岩石
不应同时保留两者。

\textbf{16.} $0.030$ 为透\omterm{def:b1:rock-forming-minerals:groups}{辉石}的\omterm{def:b1:optical-mineralogy:birefringence}{双折射率}；其高\omterm{def:b1:optical-mineralogy:relief}{突起}、
无色及接近 \ang{87} 的两组解理相符：B 为
单斜\omterm{def:b1:rock-forming-minerals:groups}{辉石}，即透\omterm{def:b1:rock-forming-minerals:groups}{辉石}。

\textbf{17.} $\qty{30000}{nm}\times0.030 = \qty{900}{nm}$：二级黄。

\textbf{18.} 厚度至 \qty{8}{\micro m} 时，石英至多显示 \qty{270}{nm}，
为一级白。

\textbf{19.} \omterm{def:b1:optical-mineralogy:twin}{聚片双晶}：D 为\omterm{def:b1:rock-forming-minerals:groups}{斜长石}。

\textbf{20.} \omterm{def:b1:rock-forming-minerals:groups}{斜长石}的成分（即其钙\omterm{def:b1:rock-forming-minerals:groups}{长石}含量）。

\textbf{21.} \omterm{def:b1:optical-mineralogy:extinction}{消光角}为 \ang{17} 的\omterm{def:b1:rock-forming-minerals:groups}{角闪石}：普通\omterm{def:b1:rock-forming-minerals:groups}{角闪石}
（\ang{15}--\ang{19}），而非\omterm{def:b1:rock-forming-minerals:groups}{辉石}（透\omterm{def:b1:rock-forming-minerals:groups}{辉石}为 \ang{38}）。

\textbf{22.} $\qty{38000}{nm}\times0.020 = \qty{760}{nm}$：二级
绿。

\textbf{23.} 该切面垂直于\omterm{def:b1:optical-mineralogy:optic-axis}{光轴}切割：其两个折射率
对沿\omterm{def:b1:optical-mineralogy:optic-axis}{光轴}传播的光相等，$\Delta = 0$，故
与均质颗粒一样保持消光。

\textbf{24.} A 为石英，B 为透\omterm{def:b1:rock-forming-minerals:groups}{辉石}，C 为黑\omterm{def:b1:rock-forming-minerals:groups}{云母}，D 为\omterm{def:b1:rock-forming-minerals:groups}{斜长石}，E 为普通\omterm{def:b1:rock-forming-minerals:groups}{角闪石}：一种
含少量石英的中性火成岩，如石英闪长岩。

\textbf{25.} 该\omterm{def:b1:field-lab-methods-1:thin-section}{薄片}厚 \textbf{\qty[text-series-to-math=true]{38}{\micro m}}。
\end{solution}
